\documentclass[aspectratio=169]{beamer}
\input{header}
\coursecode{JKSSB}

\title{Current CPU Architecture, Data Representation \& RISC/CISC}
\subtitle{Lesson 59 --- von Neumann, Addressing, Hazards, BCD, RISC vs CISC}
\author{JKSSB Computer Awareness}
\date{\today}

\begin{document}
\frame{\titlepage}

\begin{frame}{von Neumann Architecture \& the Bottleneck}
  \begin{itemize}
    \item \key{von Neumann (stored-program) architecture}: instructions and
      data are held together in the \key{same memory}.
    \item The CPU repeatedly fetches an instruction, decodes it, and executes
      it, fetching data from the same memory.
    \item A single shared path carries both instructions and data.
  \end{itemize}
  \begin{alertblock}{von Neumann bottleneck}
    The limited rate at which instructions and data can move between the CPU
    and memory on that shared path. The CPU often waits for memory, which
    caps performance.
  \end{alertblock}
\end{frame}

\begin{frame}{Addressing Modes: Immediate}
  \begin{itemize}
    \item An \key{addressing mode} is how an instruction specifies where its
      operand comes from.
    \item \key{Immediate addressing}: the operand is carried \emph{inside}
      the instruction itself.
    \item No extra memory access is needed to fetch the operand.
  \end{itemize}
  \begin{block}{Immediate in one line}
    The operand is part of the instruction, not a separate memory location.
  \end{block}
  \muted{Contrast: register addressing names a register; direct addressing names a memory address.}
\end{frame}

\begin{frame}{Cache, Hit Ratio \& Why a Hit Is Faster}
  \begin{itemize}
    \item A \key{cache} is a small, fast memory between the CPU and main
      memory that holds recently used instructions and data.
    \item A \key{hit} means the requested item is present; a \key{miss} means
      it is not.
  \end{itemize}
  \begin{block}{Cache-hit ratio}
    $\text{hit ratio} = \dfrac{\text{hits}}{\text{hits} + \text{misses}}
    = \dfrac{\text{hits}}{\text{total accesses}}$.
  \end{block}
  \begin{alertblock}{Why a hit is faster}
    Cache access time is far lower than main-memory access time, so a hit
    avoids a slow main-memory access.
  \end{alertblock}
\end{frame}

\begin{frame}{Data Hazards Even with Register Addressing}
  \begin{itemize}
    \item A \key{data hazard} is a dependency that stops an instruction from
      proceeding correctly in the expected cycle.
    \item Register operands avoid an extra memory access, but they do
      \bad{not} remove every hazard.
    \item A read-after-write dependency (an instruction needs a register that
      an earlier instruction is still writing) is a hazard even though no
      memory operand is involved.
  \end{itemize}
  \begin{block}{Key point}
    Hazards come from \emph{dependencies between instructions}, typically
    exposed by pipelining, not only from memory operands.
  \end{block}
\end{frame}

\begin{frame}{Virtual Memory vs Cache}
  \begin{columns}[T]
    \column{0.48\textwidth}
      \textbf{Cache}
      \begin{itemize}
        \item Small, fast memory close to the CPU.
        \item Holds recently used blocks.
        \item Managed mostly in hardware.
        \item Goal: lower average access time.
      \end{itemize}
    \column{0.48\textwidth}
      \textbf{Virtual memory}
      \begin{itemize}
        \item Makes memory appear larger than physical RAM.
        \item Uses secondary storage as backing store.
        \item Managed by the OS with hardware help.
        \item Goal: run large programs within limited RAM.
      \end{itemize}
  \end{columns}
  \vspace{0.2cm}
  \muted{The two solve different problems and are not the same hardware.}
\end{frame}

\begin{frame}{Instruction Cycle: Program Counter \& Instruction Register}
  \begin{itemize}
    \item \key{Program Counter (PC)}: holds the address of the \key{next}
      instruction to be fetched.
    \item \key{Instruction Register (IR)}: holds the instruction currently
      being fetched, decoded, and executed.
    \item After fetching, the PC is advanced to point at the following
      instruction.
  \end{itemize}
  \begin{alertblock}{Trap}
    The PC holds the address of the \emph{next} instruction, \bad{not} the
    current one being executed.
  \end{alertblock}
\end{frame}

\begin{frame}{Interrupts at Instruction Boundaries}
  \begin{itemize}
    \item An \key{interrupt} is a signal that makes the CPU pause the current
      program and run a handler.
    \item Interrupts are checked at \key{instruction boundaries}, during the
      program run.
    \item They are \good{not} delayed until the whole program finishes.
  \end{itemize}
  \begin{alertblock}{Trap}
    ``Interrupts are processed only after the entire program completes'' is
    false. Handling happens mid-program, between instructions.
  \end{alertblock}
\end{frame}

\begin{frame}{Instruction Cycle Counts Vary}
  \begin{itemize}
    \item Different instructions do different amounts of work.
    \item A simple register operation may complete in one clock cycle.
    \item A divide, a multi-step operation, or an instruction that accesses
      slow memory can take many cycles.
  \end{itemize}
  \begin{alertblock}{Trap}
    ``All machine instructions take the same number of clock cycles'' is
    false. The count \bad{varies} by instruction.
  \end{alertblock}
\end{frame}

\begin{frame}{Two's-Complement \& Overflow Detection}
  \begin{itemize}
    \item In $n$-bit two's complement the most significant bit is the
      \key{sign} bit; negative values are formed by inverting and adding one.
    \item \key{Overflow} means the true result cannot be represented in the
      available bits.
  \end{itemize}
  \begin{block}{Overflow test}
    Overflow occurs when two operands of the \key{same sign} produce a result
    of the \key{opposite sign}.
  \end{block}
  \begin{alertblock}{Trap}
    Overflow is \bad{not} the same as a carry out of the most significant
    bit. A carry can occur without any overflow.
  \end{alertblock}
\end{frame}

\begin{frame}{Zero-Extension vs Sign-Extension}
  \begin{itemize}
    \item \key{Extension} widens a value to more bits while keeping its
      value.
    \item \key{Zero-extension}: fill the new leftmost bits with 0 --- used
      for \key{unsigned} values.
    \item \key{Sign-extension}: copy the sign bit into the new leftmost bits
      --- used for \key{signed} values.
  \end{itemize}
  \begin{block}{Remember the pair}
    Unsigned $\rightarrow$ zero-extension; signed $\rightarrow$
    sign-extension. Swapping them is the common trap.
  \end{block}
\end{frame}

\begin{frame}{BCD Encoding}
  \begin{itemize}
    \item \key{BCD} = Binary-Coded Decimal.
    \item Each decimal digit is stored in its own \key{4-bit} code.
    \item Decimal $59$ becomes $0101\ 1001$ (5 then 9).
  \end{itemize}
  \begin{block}{Contrast with plain binary}
    Plain binary of $59$ is $111011_2$; BCD keeps one 4-bit code per decimal
    digit.
  \end{block}
  \begin{alertblock}{Trap}
    BCD does \bad{not} pack the whole decimal number as ordinary binary.
  \end{alertblock}
\end{frame}

\begin{frame}{Byte-Addressability \& Effective Address}
  \begin{itemize}
    \item \key{Byte-addressable memory}: every individual byte has its own
      unique address; a word is a group of bytes at consecutive addresses.
    \item The \key{effective address} is the actual location an instruction
      uses for its operand; it may be computed at run time.
    \item Instruction length (in bytes) and effective address are
      \key{different} quantities.
  \end{itemize}
  \begin{block}{Key point}
    A long instruction does not mean a far address; the operand address is
    determined by the addressing mode, not by the instruction's size.
  \end{block}
\end{frame}

\begin{frame}{RISC vs CISC}
  \begin{center}
    \begin{tabular}{lll}
      \toprule
      \textbf{Aspect} & \textbf{RISC} & \textbf{CISC} \\
      \midrule
      Instruction set & Small, simple & Large, complex \\
      Instruction length & Fixed & Variable \\
      Registers & Many & Fewer \\
      Memory access & Load/store only & Memory operands allowed \\
      Cycles per instruction & Typically one (pipelined) & Often multiple \\
      Control emphasis & Hardwired, compiler-driven & Microcoded, hardware-driven \\
      \bottomrule
    \end{tabular}
  \end{center}
  \vspace{0.2cm}
  \muted{RISC = Reduced Instruction Set Computer; CISC = Complex Instruction Set Computer.}
\end{frame}

\begin{frame}{Lesson 59: Recall}
  \begin{itemize}
    \item von Neumann: instructions and data share one memory; the shared
      path causes the von Neumann bottleneck.
    \item Immediate addressing keeps the operand inside the instruction.
    \item Cache-hit ratio $=$ hits $/$ total accesses; a hit is faster than
      main memory.
    \item Hazards come from dependencies, even with register addressing.
    \item The PC holds the address of the \key{next} instruction; the IR
      holds the current instruction.
    \item Interrupts are handled at instruction boundaries, mid-program.
    \item Clock-cycle counts vary by instruction.
    \item Overflow $=$ same-sign operands give an opposite-sign result.
    \item Zero-extension for unsigned; sign-extension for signed.
    \item BCD uses one 4-bit code per decimal digit.
    \item RISC $=$ simple fixed-length, many registers, load/store.
  \end{itemize}
\end{frame}
\end{document}
